\section{Discussion and limitations}
\label{sec:limitations}
The answer depends on what investment buys. If it buys independent, fully controllable task-score improvements, both snapshots show a large Borda-versus-mean task-budget gap under the primary setting. If it buys only a fraction of available headroom, first place becomes infeasible for most targets under several rules. If it buys correlated improvements, unequal task costs or changes in the rival population, the optimization problem itself changes. This study measures the first two settings and does not identify the third from score tables.

A unit task cost is useful as a measure of intervention breadth, not as a physical resource. MMLU subjects and broad Open LLM tasks have different sizes, formats and difficulty. Neither a fraction of 57 subjects nor a fraction of six tasks can be translated into tokens, dollars or developer effort without further measurements. Full score access is particularly optimistic: some selected scores may be unattainable by training, and attaining them jointly may incur interference. Even an exact feasible subset is a certificate only within the declared score-space intervention.

Higher task budgets alone are not a sufficient benchmark-design objective. Borda's rankings depend on which rivals are present and discard score margins; minimum aggregation can focus attention on a single low threshold; and changing normalization can alter the intervention menu as well as the ranking. We assess no trade-off between strategic budget, statistical reliability, intended task utility and genuine model quality. A designer should specify those goals and plausible gain/cost models before treating any rule as preferable.

The empirical scope is two fixed public snapshots with no representative sampling of models or tasks. Shared model ancestry, duplicated evaluations, correlated subjects and task-level score estimation error limit generalization. We address one deterministic deduplication scheme and several composition perturbations, but cannot infer score uncertainty from model-to-model dispersion. No actual contamination, targeted training or live leaderboard manipulation is performed. Practical safeguards against contamination require evidence beyond these score edits. Release of counterfactual budgets may also reveal incentive weaknesses; our purpose is to make evaluation assumptions inspectable and support defensive benchmark design.

\section{Conclusion}
Minimum task-count investment can be computed exactly under an explicit score-improvement menu. On the two measured normalized snapshots, primary common challengers need substantially more selected tasks for Borda than for arithmetic mean, while conventional robust cardinal rules do not imply a uniformly high strategic budget. Bounded gains, task composition, normalization and the definition of first place qualify that comparison. The evidence supports conditional score-space conclusions and a reproducible way to test them; actual training costs and attainable gains remain unmeasured.
